% !TeX encoding = UTF-8
\documentclass[variant=docente,cover=institutional,mode=screen]{ua-engineering-v6-article}
% !TeX encoding = UTF-8
\UASetUniversity{Universidad Austral}
\UASetFaculty{Facultad de Ingeniería}
\UASetDepartment{Departamento de Computación}
\UASetCampus{Pilar}
\UASetCareer{Ingeniería Informática}
\UASetCommission{Comisión A}
\UASetProfessor{Equipo docente}
\UASetSemester{Segundo cuatrimestre 2026}
\UASetCourse{Sistemas Distribuidos}
\UASetDate{2026}

\UASetSemester{Segundo cuatrimestre 2026}
\UASetDate{2026}
\usetikzlibrary{positioning,arrows.meta,calc,fit,matrix,decorations.pathreplacing}

\UASetClassNumber{06}
\UASetClassTitle{Storage distribuido y almacenamiento permanente}
\UASetDocKind{Guía docente magistral y dossier técnico}
\title{Clase 06 --- Guía docente: almacenamiento permanente, acceso concurrente y rendimiento bajo fallas}
\author{\UAProfessor}
\date{\UADate}

\begin{document}
\setlength{\emergencystretch}{2em}
\maketitle
\tableofcontents
\clearpage

\section{Parte A: guion operativo para 180 minutos}

\subsection{Objetivo pedagógico profundo}

Esta clase inaugura el bloque de storage y debe producir dos efectos simultáneos:

\begin{enumerate}
\item recuperar de manera ágil el lenguaje de incertidumbre, reintentos, replicación, consenso y transacciones;
\item mostrar que las garantías lógicas existen porque mecanismos locales convierten decisiones en bytes recuperables, versiones visibles, rutas físicas y trabajo de mantenimiento controlado.
\end{enumerate}

El riesgo principal es presentar WAL, MVCC, \texttt{EXPLAIN}, LSM, filtros de Bloom, compactación y stalls como una lista de términos. La secuencia de enseñanza debe ser:

\[
\begin{aligned}
\text{contexto} &\rightarrow \text{problema observable} \rightarrow \text{intuición} \\
&\rightarrow \text{definición} \rightarrow \text{timeline o plan} \\
&\rightarrow \text{falla} \rightarrow \text{evidencia} \rightarrow \text{decisión}.
\end{aligned}
\]

\begin{uakeybox}
\textbf{Resultado cognitivo central.} Al finalizar, un estudiante debe poder mirar un COMMIT, un snapshot, un plan de ejecución o un dashboard de compactación y explicar qué garantía se intenta sostener, ante qué falla, con qué costo y mediante qué evidencia.
\end{uakeybox}

La clase no es una capacitación de PostgreSQL o RocksDB. Ambos son casos reales para mostrar que los principios se materializan en motores concretos.

\subsection{Resultados de aprendizaje observables}

Al terminar, el estudiante debería poder:

\begin{itemize}
\item distinguir capa distribuida de motor de almacenamiento local;
\item explicar todas las siglas antes de utilizarlas;
\item describir qué representa el ID 9F2A y por qué debe ser durable;
\item diferenciar operación aceptada, escrita, committed, durable, aplicada, visible y replicada;
\item explicar las dos reglas de WAL y completar el timeline de CampusShop;
\item responder correctamente si el flush durable ``empieza'' en el \texttt{WAL UPDATE};
\item resolver cuatro puntos de crash y justificar ``todas y solo las confirmadas'';
\item explicar group commit y checkpoint sin confundirlos con pérdida de garantía o backup;
\item definir snapshot, visibilidad, vista coherente y serializabilidad;
\item explicar por qué MVCC no elimina write skew ni el costo de limpiar versiones;
\item leer \texttt{EXPLAIN} de abajo hacia arriba y justificar un índice compuesto;
\item describir la ruta de escritura y lectura de una LSM-tree;
\item explicar filtro de Bloom, merge ordenado, tombstone, backlog y write stall;
\item interpretar conjuntamente ingreso, capacidad de compactación, archivos L0 y p99;
\item comparar una solución para CampusShop y otra para telemetría del campus.
\end{itemize}

\section{Arquitectura didáctica v9}

La clase se organiza en bloques de exposición breve, predicción, representación visual, actividad y cierre. El objetivo no es recorrer todos los términos linealmente, sino alternar explicación y evidencia para sostener la comprensión durante las tres horas.

\subsection{Arquitectura temporal de 180 minutos}

\begin{center}
\begin{tabular}{p{0.12\textwidth}p{0.24\textwidth}p{0.53\textwidth}}
\toprule
Tiempo & Bloque & Propósito y evidencia \\
\midrule
0--10 & Apertura & Compra confirmada y corte de energía; predicción individual y en parejas. \\
10--25 & Contexto & Storage distribuido, capa distribuida, motor local y objetos físicos. \\
25--42 & Vocabulario & WAL, MVCC, LSM y términos de durabilidad, recovery, concurrencia y SQL. \\
42--52 & Repaso 1--5 & Cinco frases, prediagnóstico intuitivo y mapa conceptual. \\
52--82 & WAL & Reglas, carriles, timeline, flush del prefijo, crashes y recovery. \\
82--92 & Actividad + pausa & Equipos completan timeline; pausa breve mientras se comparan resultados. \\
92--119 & MVCC & Reporte largo, snapshot, versiones, VACUUM y write skew. \\
119--145 & Índices y EXPLAIN & Consulta, planes sin/con índice, costo de escritura y mini demo. \\
145--169 & LSM & Write path, merge, Bloom, read path, backlog, compaction y stall. \\
169--178 & Integración & CampusShop vs telemetría; defensa de 90 segundos. \\
178--180 & Cierre & Recuperación del diagnóstico y exit ticket. \\
\bottomrule
\end{tabular}
\end{center}

No conviene hablar más de 15 minutos sin cambiar el modo de actividad. Alternar explicación, predicción, dibujo, cálculo, demo y discusión.

\subsection{Selección de diapositivas: núcleo, actividad y respaldo}

La presentación tiene construcciones progresivas; no todas las páginas físicas son conceptos distintos. En una observación docente conviene declarar que las etiquetas cumplen esta función:

\begin{description}
\item[NÚCLEO.] Contenido indispensable para lograr el resultado de aprendizaje.
\item[ACTIVIDAD.] Producción observable del estudiante.
\item[EXTENSIÓN.] Profundización si el ritmo y las preguntas lo permiten.
\item[RESPALDO.] Material para responder preguntas o ampliar, no para exposición lineal obligatoria.
\end{description}

Si el tiempo se reduce, preserve:

\begin{enumerate}
\item contexto y vocabulario;
\item timeline WAL con criterio de ACK;
\item timeline MVCC y write skew;
\item transición MVCC--árbol B y comparación de planes;
\item LSM, Bloom, backlog y stall;
\item actividad final y exit ticket.
\end{enumerate}

La correspondencia con la versión actual de las diapositivas es:

\begin{center}
\small
\begin{tabular}{p{0.16\textwidth}p{0.29\textwidth}p{0.43\textwidth}}
\toprule
Frames fuente & Bloque & Producción esperada \\
\midrule
1--18 & Contexto, vocabulario, diagnóstico & Definir promesa, capas, acrónimos y evidencias \\
19--30 & WAL y recuperación & Timeline, prefijo durable, ACK y matriz de crashes \\
31--39 & MVCC & Versiones visibles, bloat y write skew \\
40--49 & Árbol B, índice y \texttt{EXPLAIN} & Ruta física, plan sin/con índice y costo de mantenerlo \\
50--63 & LSM, Bloom, compactación y stall & Write/read path, backlog y amplificaciones \\
64--69 & Integración & Decisión CampusShop/telemetría, repaso y exit ticket \\
\bottomrule
\end{tabular}
\end{center}

\subsection{Preparación previa}

\subsubsection*{Cuarenta y ocho horas antes}

\begin{itemize}
\item Compilar la presentación, apuntes, guía, soluciones y ficha de aula.
\item Imprimir o dibujar un timeline WAL de cinco carriles y uno MVCC de cuatro carriles.
\item Ejecutar la demo SQL de \texttt{EXPLAIN}; guardar también la salida para usar si falla el laboratorio.
\item Preparar el dashboard didáctico de compactación con ingreso 90 MB/s, capacidad 60 MB/s y backlog 18 GB.
\item Ensayar las cinco frases de repaso, con un máximo de 40 segundos por clase previa.
\item Revisar la semántica de \texttt{fsync} y las limitaciones del hardware; evitar prometer más de lo que el modelo declara.
\item Releer la definición de filtro de Bloom y practicar la diferencia entre falso positivo y falso negativo.
\item Resolver la guía completa sin consultar las soluciones y marcar cualquier respuesta que todavía necesite estudio.
\end{itemize}

\subsubsection*{Treinta minutos antes}

\begin{itemize}
\item Escribir en el pizarrón la pregunta central.
\item Dibujar las capas: operación lógica, capa distribuida, motor local y estado permanente.
\item Dejar preparados los carriles WAL y MVCC sin eventos.
\item Abrir la consulta SQL y las dos salidas de plan.
\item Escribir la igualdad: ingreso $-$ capacidad $=$ crecimiento del backlog.
\item Verificar que las etiquetas de la presentación se leen desde el fondo del aula.
\end{itemize}

\subsubsection*{Cinco minutos antes}

\begin{itemize}
\item Probar el proyector y el tamaño del código.
\item Tener un cronómetro visible.
\item Repartir la ficha o habilitar una versión digital.
\item Preparar dos tarjetas: \texttt{COMMIT DURABLE} y \texttt{ACK}.
\end{itemize}

\subsection{Plan de pizarrón}

Dividir el pizarrón en tres zonas y no borrar las dos primeras hasta el cierre.

\begin{enumerate}
\item \textbf{Izquierda --- promesa y vocabulario.} CampusShop, 9F2A, garantía ``todas y solo las confirmadas'', definición de durable.
\item \textbf{Centro --- ejecuciones.} Timeline WAL arriba y timeline MVCC abajo.
\item \textbf{Derecha --- evidencia.} Plan sin/con índice y cadena ingreso--backlog--stall.
\end{enumerate}

En el extremo superior mantener visible:

\[
\text{problema} \rightarrow \text{mecanismo} \rightarrow
\text{garantía} \rightarrow \text{costo} \rightarrow \text{evidencia}.
\]

\subsection{Demostraciones}

Las demostraciones no son intermedios decorativos: cada una debe comprobar una afirmación concreta y terminar con una evidencia que los estudiantes puedan explicar.

\begin{enumerate}
\item \textbf{Recovery conceptual.} Entregar tarjetas \texttt{BEGIN}, \texttt{WAL UPDATE}, \texttt{COMMIT}, \texttt{ACK}, \texttt{CRASH} y \texttt{REDO}. Mover el crash antes y después de la frontera durable y pedir que indiquen qué debe aparecer tras recovery.
\item \textbf{Plan de ejecución.} Mostrar la misma consulta con \texttt{EXPLAIN ANALYZE} antes y después de crear el índice compuesto. Comparar tipo de scan, filas, sort, buffers y costo agregado a escrituras.
\item \textbf{MVCC humano.} Un estudiante representa un reporte con snapshot antiguo; otro confirma una nueva versión. Preguntar qué versión ve cada lector y cuándo la anterior puede reclamarse.
\item \textbf{Compactación manual.} Fusionar dos listas ordenadas y un tombstone, contando bytes leídos y reescritos. Relacionar el resultado con amplificación, backlog y write stall.
\end{enumerate}

Para cada demostración, cierre con cuatro preguntas: qué problema resuelve, qué garantiza, qué no garantiza y qué costo desplaza.

\subsection{Guion de apertura, 0--10 minutos}

Diga:

\begin{quote}
``Una estudiante compra la última entrada disponible. La aplicación muestra `compra confirmada'. Un segundo después se corta la energía. No quiero todavía una sigla: quiero saber qué debe seguir siendo verdadero cuando el sistema reinicie.''
\end{quote}

Dé 60 segundos individuales y 90 segundos en parejas. Recoja respuestas sin corregir de inmediato. Es probable que aparezcan ``RAM'', ``disco'', ``réplica'', ``base'' y ``log''. Pregunte:

\begin{itemize}
\item ¿disco de cuál nodo?;
\item ¿una réplica en RAM sobrevive a un corte común?;
\item ¿qué significa exactamente que ``la base confirmó''?;
\item ¿cómo reconoce el mismo intento si la respuesta se perdió?
\end{itemize}

Cierre el bloque:

\begin{quote}
``Hoy vamos a convertir esas palabras en fronteras comprobables.''
\end{quote}

\subsection{Contexto del storage, 10--25 minutos}

\subsubsection*{Paso 1: storage distribuido}

Explique:

\begin{quote}
``La capa distribuida decide dónde vive el dato, quién tiene autoridad y qué réplicas deben participar. El motor local decide cómo esa decisión sobrevive un crash y cómo vuelve a ser consultable.''
\end{quote}

Pregunte qué parte de Raft pertenece a la capa distribuida y qué estado debe persistir cada nodo. La respuesta esperada es que el consenso ordena y decide, pero votos, término y entradas deben sobrevivir según el protocolo.

\subsubsection*{Paso 2: motor de almacenamiento}

No utilice todavía nombres de productos. Dibuje gestor transaccional, planificador, memoria intermedia, WAL/recovery, estructuras y medio persistente.

Frase sugerida:

\begin{quote}
``El motor es el traductor entre `confirmar compra' y bytes, versiones, índices y reglas de recuperación.''
\end{quote}

\subsubsection*{Paso 3: objetos del caso}

Defina 9F2A antes de usarlo:

\begin{quote}
``9F2A no es una dirección ni un LSN. Es la identidad estable de esta compra lógica. Si la red hace repetir el request, la identidad no cambia.''
\end{quote}

Ubique fila, tabla, página P17, índice I5 y registro WAL. Pregunte cuál es lógico y cuál físico.

\subsection{Vocabulario, 25--42 minutos}

No lea tablas completas. Use el patrón:

\begin{enumerate}
\item pronunciar el término y expandir la sigla;
\item traducirlo al español;
\item dar una frase intuitiva;
\item dar una definición técnica;
\item marcar un límite.
\end{enumerate}

Ejemplo para WAL:

\begin{quote}
``WAL significa Write-Ahead Logging, registro anticipado de escritura. Intuitivamente es el cuaderno de evidencia del motor. Técnicamente impone que la evidencia de recuperación llegue al log durable antes de depender de la página, y que COMMIT sea durable antes del ACK. No reemplaza replicación ni idempotencia.''
\end{quote}

Para durable, insistir:

\begin{quote}
``Durable no significa `muy guardado'. Significa reconstruible después de una falla específica.''
\end{quote}

Haga que un estudiante defina \emph{flush común} y otro \emph{flush durable}. La distinción esperada es que el primero mueve bytes; el segundo afirma que ya cruzaron la frontera del modelo.

\subsection{Repaso y prediagnóstico, 42--52 minutos}

Use exactamente una frase por clase previa:

\begin{itemize}
\item Clase 1: incertidumbre requiere evidencia persistida.
\item Clase 2: retry seguro requiere identidad y resultado durables.
\item Clase 3: muchas copias no dicen cuáles son durables o actuales.
\item Clase 4: consenso decide qué recordar; storage hace que se recuerde.
\item Clase 5: ACID define el contrato; WAL y MVCC implementan parte.
\end{itemize}

En el prediagnóstico no permita respuestas por siglas. Pregunte primero ``¿qué comportamiento sería correcto?''. Registre votos para volver al final.

\subsection{WAL y recuperación, 52--82 minutos}

\subsubsection*{A. Declarar la garantía}

Antes de dibujar WAL, escriba:

\begin{quote}
``Después de recovery: todas y solo las transacciones confirmadas.''
\end{quote}

Pregunte qué significa ``todas'' y qué significa ``solo''. Esta frase permite justificar los tres resultados de crash sin memorizar casos.

\subsubsection*{B. Jerarquía y frontera durable}

Dibuje aplicación, buffer, page cache, cache de dispositivo y medio persistente. Pregunte dónde puede terminar \texttt{write()}.

No diga que \texttt{fsync} es una magia universal. Diga:

\begin{quote}
``La garantía depende de la semántica real de todo el camino y del tipo de falla que declaramos.''
\end{quote}

\subsubsection*{C. Dos reglas}

Para cada regla, primero muestre el orden correcto; después inviértalo.

\begin{itemize}
\item Log antes que página: evita páginas adelantadas sin evidencia suficiente.
\item COMMIT antes que ACK: evita promesas al cliente que recovery no puede sostener.
\end{itemize}

\subsubsection*{D. Construcción del timeline}

Pida que ubiquen cada evento en el carril correcto. No muestre de inmediato la solución completa.

La frase clave del paso 6 es:

\begin{quote}
``Ahora hacemos flush del WAL hasta el LSN del COMMIT. No empezamos el flush en el paso del UPDATE: el UPDATE se agregó antes; el flush de ahora vuelve durable el prefijo requerido.''
\end{quote}

Escriba:

\[
\texttt{flushLSN} \geq \texttt{commitLSN}(T).
\]

Pregunte por qué eso incluye los updates anteriores. Respuesta: porque WAL es secuencial y la frontera avanza sobre un prefijo.

\subsubsection*{E. Mover el crash}

Mueva una tarjeta CRASH a cuatro posiciones. Para cada una, pida:

\begin{enumerate}
\item qué observó el cliente;
\item qué evidencia sobrevivió;
\item qué debe hacer recovery;
\item si el retry con 9F2A crea otro efecto.
\end{enumerate}

No simplifique el caso ``COMMIT agregado pero no flushed'': para el cliente el resultado es incierto. Si algunos bytes sobreviven, la transacción puede reaparecer aunque el ACK no haya llegado.

Anticipe una pregunta frecuente: ``si las páginas no se escribieron, ¿de dónde sale toda la data?''. Responda sobre el mismo dibujo:

\begin{quote}
``Recovery no reconstruye la base desde la palabra COMMIT. Parte de las páginas persistentes anteriores y aplica los registros WAL de modificación que faltan. Los UPDATE del WAL contienen cómo rehacer; COMMIT indica que ese conjunto de cambios pertenece a la historia confirmada.''
\end{quote}

Escriba:

\[
\text{páginas base} + \operatorname{REDO}(\text{WAL confirmado pendiente})
= \text{estado recuperado}.
\]

Aclare el límite: pérdida completa del dispositivo requiere backup base más WAL archivado o una réplica adecuada.

\subsubsection*{F. Group commit y checkpoint}

Group commit se explica con tres tarjetas de transacción y un solo \texttt{fsync}. Pregunte si la garantía cambia. Respuesta: no, mientras cada ACK espere el flush que contiene su COMMIT.

Checkpoint se introduce como respuesta a ``¿por qué recovery no relee desde el origen?''. Diferéncielo de backup y réplica.

\subsection{Actividad WAL y pausa, 82--92 minutos}

En grupos de tres, entregar un timeline vacío. Cada grupo debe marcar:

\begin{itemize}
\item el prefijo durable;
\item el punto más temprano de ACK legítimo;
\item un crash que exige REDO;
\item un crash con resultado desconocido para el cliente;
\item el papel de 9F2A.
\end{itemize}

Mientras comparan, hacer una pausa breve. Conservar dos producciones diferentes para discutir.

\subsection{MVCC, 92--119 minutos}

\subsubsection*{A. Problema antes del mecanismo}

Diga:

\begin{quote}
``Un reporte tarda 40 segundos y entran 900 compras. ¿Bloqueamos checkout, mezclamos momentos o conservamos historia?''
\end{quote}

\subsubsection*{B. Definir snapshot correctamente}

No usar solamente ``foto''. Diga:

\begin{quote}
``Snapshot es una regla de visibilidad: determina qué transacciones confirmadas y qué versiones pertenecen a la historia que puede observar esta lectura.''
\end{quote}

Marque el límite: no es copia completa ni garantía automática de serializabilidad.

\subsubsection*{C. Timeline humano}

Un estudiante representa $T_R$, otro $T_W$, y dos tarjetas representan $v_1$ y $v_2$. Mantenga $v_1$ visible para el lector antiguo después del COMMIT de $v_2$. Pregunte qué lee un lector nuevo.

\subsubsection*{D. Costo de conservar historia}

Introduzca VACUUM solo después de que entiendan por qué no se puede borrar $v_1$. Frase:

\begin{quote}
``Una lectura larga conserva el pasado y puede hacer más caro el presente.''
\end{quote}

\subsubsection*{E. Write skew}

Recupere el ejemplo de dos guardias. Pregunte qué conflicto no detectó el sistema. Concluya:

\begin{quote}
``MVCC es un mecanismo de versiones; serializabilidad es una garantía sobre ejecuciones.''
\end{quote}

\subsection{Índices y EXPLAIN, 119--145 minutos}

\subsubsection*{A. Transición suave desde MVCC}

Antes de mostrar la consulta, escriba dos preguntas:

\begin{itemize}
\item MVCC: ``¿qué versión puede observar esta transacción?''
\item Índice: ``¿cómo encuentra el motor las versiones candidatas sin revisar todas las páginas?''
\end{itemize}

Diga:

\begin{quote}
``El índice no decide qué versión es correcta. Reduce el espacio de búsqueda. MVCC aplica después la regla de visibilidad.''
\end{quote}

Esta transición evita que el árbol B aparezca como un tema aislado.

\subsubsection*{B. Partir de la pregunta de producto}

Mostrar la consulta completa y subrayar igualdad por estudiante, igualdad por estado, orden descendente por fecha y \texttt{LIMIT 20}. Antes de mostrar el índice, pedir propuestas.

Construya luego el árbol B en dos pasos: primero raíz, páginas internas y hojas; después la ruta de CampusShop. Defina \emph{fan-out} como la cantidad de claves y referencias que caben en una página. No presente el árbol como un dibujo abstracto: cada nodo representa una página del índice y las hojas contienen claves ordenadas más referencias.

\subsubsection*{C. EXPLAIN no es el índice}

Diga:

\begin{quote}
``El índice es una estructura. EXPLAIN es evidencia sobre la ruta que el planificador eligió. EXPLAIN ANALYZE ejecuta y mide.''
\end{quote}

\subsubsection*{D. Leer de abajo hacia arriba}

Pida que un estudiante narre:

\[
\text{Seq Scan} \rightarrow \text{Filter} \rightarrow \text{Sort} \rightarrow \text{Limit}.
\]

Después compare con Index Scan. No comenzar por el tiempo; comenzar por filas, buffers, operador y orden.

\subsubsection*{E. Costo del índice}

Por cada índice secundario, agregar una tarjeta al camino de \texttt{INSERT}. Preguntar qué aumenta: WAL, páginas, caché, espacio y mantenimiento.

\subsubsection*{F. Demo y fallback}

Si PostgreSQL está disponible, ejecutar la consulta y \texttt{EXPLAIN (ANALYZE, BUFFERS)} antes y después del índice. Si no, usar las salidas impresas. La clase no debe depender de una instalación en vivo.

\subsection{LSM, Bloom, compactación y stalls, 145--169 minutos}

\subsubsection*{A. LSM antes de compaction}

Definir la sigla y dibujar el write path:

\[
WAL \rightarrow memtable \rightarrow memtable\ inmutable
\rightarrow flush \rightarrow SSTable\ L0.
\]

Aclarar que LSM no elimina el costo: desplaza organización fuera del camino crítico.

\subsubsection*{B. Merge ordenado}

Use dos listas ya ordenadas en tarjetas. Fusiónelas comparando las cabeceras. Explique que no se ordena todo desde cero.

\subsubsection*{C. Tombstone}

Introduzca una tarjeta ``eliminar c''. Pregunte por qué no puede descartarse si una SSTable vieja aún contiene $c$. La respuesta deseada es evitar resurrección.

\subsubsection*{D. Filtro de Bloom}

Diga:

\begin{quote}
``Bloom no busca el valor. Es un resumen probabilístico que permite demostrar algunas ausencias. Si dice `no está', omitimos la SSTable. Si dice `tal vez', todavía debemos consultar.''
\end{quote}

No usar ``archivo imposible''. Usar ``SSTable que puede omitirse para esta búsqueda''.

Pregunte:

\begin{itemize}
\item ¿puede haber falso positivo? Sí.
\item ¿puede haber falso negativo? No, si el filtro está correctamente construido y consultado.
\end{itemize}

\subsubsection*{E. Dashboard y stall}

Escriba:

\[
90-60=30\ \text{MB/s de backlog adicional}.
\]

Pida calcular 5 minutos hasta slowdown y 10 hasta stop desde 18 GB, usando unidades decimales. Después pregunte si el stall es error. Respuesta: es protección; su frecuencia y duración pueden revelar incapacidad.

\subsection{Integración y cierre, 169--180 minutos}

Equipo A defiende CampusShop; equipo B, telemetría. Cada presentación debe contener:

\begin{enumerate}
\item problema;
\item falla;
\item garantía;
\item mecanismo;
\item costo;
\item métrica.
\end{enumerate}

Vuelva al prediagnóstico y cambie las respuestas con evidencia. Recoja el exit ticket de cuatro preguntas.

Frase de cierre:

\begin{quote}
``El storage no es una caja debajo del sistema. Es donde las promesas se convierten en bytes, esperas, recovery y riesgos medibles.''
\end{quote}

\section{Parte B: dossier técnico de preparación del docente}

\subsection{Modelo mental completo}

Para responder preguntas exigentes, mantenga esta separación:

\begin{enumerate}
\item \textbf{Semántica de negocio}: compra única, stock correcto, reporte coherente.
\item \textbf{Protocolo distribuido}: autoridad, replicación, quorum y orden.
\item \textbf{Motor local}: WAL, recovery, versiones, índices y archivos.
\item \textbf{Medio físico}: buffers, sistema operativo, controlador y dispositivo.
\item \textbf{Operación}: métricas, mantenimiento, RPO, RTO y SLO.
\end{enumerate}

Una respuesta sólida indica en qué nivel vive cada garantía y qué supuestos conecta entre niveles.

\subsection{Durabilidad: precisiones que debe dominar}

\subsubsection*{Aceptado, escrito, committed, durable, aplicado, visible y replicado}

\begin{description}
\item[Aceptado.] El sistema recibió la intención; puede no haber ejecutado nada.
\item[Escrito.] Bytes fueron copiados a alguna capa; la capa puede ser volátil.
\item[Committed.] El motor decidió confirmar la transacción según su protocolo.
\item[Durable.] La evidencia requerida sobrevivirá a la falla declarada.
\item[Aplicado.] El efecto está materializado en la estructura principal o estado derivado.
\item[Visible.] Una lectura bajo una regla concreta puede observarlo.
\item[Replicado.] Existe en más de un nodo; no indica por sí solo persistencia ni frescura.
\end{description}

Estas palabras no son intercambiables. Un COMMIT puede ser durable y todavía no estar aplicado en todas las páginas. Un dato puede estar replicado en RAM y no ser durable frente a un corte común.

\subsubsection*{Qué garantiza realmente fsync}

\texttt{fsync} solicita que el sistema operativo sincronice los datos, pero la garantía real depende de:

\begin{itemize}
\item sistema de archivos y orden de escrituras;
\item barreras y cachés del controlador;
\item cache protegida o no protegida del dispositivo;
\item configuración del motor;
\item tipo de crash considerado.
\end{itemize}

No es necesario profundizar en todos los detalles durante la clase, pero debe evitar decir ``fsync garantiza absolutamente cualquier pérdida de energía''.

\subsection{WAL en profundidad}

\subsubsection*{Por qué un log secuencial}

Persistir un log secuencial suele ser más eficiente que sincronizar muchas páginas dispersas en cada transacción. El WAL crea una representación de recuperación ordenada y permite que las páginas se escriban después.

\subsubsection*{Dos órdenes independientes}

\begin{enumerate}
\item \textbf{Update record antes que data page}: evita que la página adelante al log.
\item \textbf{Commit record durable antes que ACK}: evita que la respuesta adelante a la recuperación.
\end{enumerate}

El primer orden puede expresarse con el LSN de la página: antes de escribir una página cuyo \texttt{pageLSN}=x, el WAL debe estar durable al menos hasta $x$.

\subsubsection*{Flush del prefijo}

La pregunta ``¿el flush comienza en el WAL UPDATE?'' debe responderse así:

\begin{quote}
No. El append del update ocurre antes. El flush se solicita en un momento posterior y avanza la frontera durable hasta el LSN requerido, normalmente el COMMIT de la transacción. Ese prefijo incluye updates anteriores no persistidos.
\end{quote}

Si \texttt{flushLSN}=150 y el nuevo COMMIT está en 180, se sincroniza el tramo necesario para alcanzar 180, no se ``reinicia'' desde el primer update.

\subsubsection*{Resultado después de ACK perdido}

Si el COMMIT ya era durable y el ACK se pierde, la transacción debe permanecer confirmada. El cliente no lo sabe y reintenta. La clave 9F2A permite devolver el mismo resultado sin repetir el efecto.

WAL no ofrece exactamente-una-vez a nivel de API. Ofrece recuperación de la historia del motor. Exactamente-una-vez lógica requiere identidad y atomicidad de negocio.

\subsubsection*{REDO y UNDO}

Un diseño STEAL puede escribir páginas con cambios no confirmados y puede necesitar UNDO. Un diseño NO-STEAL evita que páginas con cambios no confirmados lleguen al medio, reduciendo UNDO pero exigiendo más memoria. FORCE escribe páginas confirmadas antes de commit; NO-FORCE permite páginas atrasadas y necesita REDO.

Muchos motores prácticos combinan STEAL + NO-FORCE con WAL, aunque los detalles varían. En clase basta con explicar que el log debe contener información suficiente para recuperar una historia válida.

\subsubsection*{Crash durante recovery}

Recovery debe ser reiniciable. \texttt{pageLSN}, checksums, registros idempotentes y fases de análisis/redo/undo son mecanismos posibles. No afirmar que un \texttt{UPDATE} SQL se vuelve a ejecutar literalmente; se reaplican registros físicos o lógicos del motor según su diseño.

\subsubsection*{Checkpoint}

Un checkpoint no necesariamente significa ``todas las páginas están perfectas en un instante''. Dependiendo del motor, registra un punto y coordina qué WAL debe conservarse y desde dónde comenzar análisis. La simplificación de la clase es útil, pero ante una pregunta avanzada aclare que existen checkpoints fuzzy y diseños ARIES-like.

\subsubsection*{Replicación y WAL}

Un sistema distribuido puede replicar registros WAL, comandos lógicos o páginas. Las preguntas son:

\begin{itemize}
\item ¿qué se envía a las réplicas?;
\item ¿cuándo se considera recibido?;
\item ¿se espera memoria o persistencia remota?;
\item ¿qué quorum habilita el ACK?;
\item ¿cómo se evita que una réplica adelante o contradiga el orden acordado?
\end{itemize}

WAL local y consenso responden problemas relacionados pero diferentes.

\subsection{MVCC en profundidad}

\subsubsection*{Qué es una versión}

Una versión suele contener el valor y metadatos de creación/eliminación. PostgreSQL utiliza identificadores de transacción; otros motores usan timestamps o secuencias. La visibilidad no depende solo del valor, sino del estado de las transacciones creadoras y reemplazantes.

\subsubsection*{Snapshot por statement o por transacción}

En Read Committed, distintas sentencias de una misma transacción pueden tomar snapshots distintos. En Repeatable Read/Snapshot Isolation suelen compartir una vista más estable. Por eso no debe afirmar que ``MVCC siempre hace que toda transacción vea exactamente la misma foto''.

\subsubsection*{Locks siguen existiendo}

MVCC reduce bloqueos entre lectores y escritores, pero no elimina:

\begin{itemize}
\item locks de escritura sobre la misma fila;
\item locks de DDL;
\item conflictos de claves únicas;
\item coordinación para serializabilidad;
\item bloqueos internos y latches.
\end{itemize}

\subsubsection*{Vacuum, versiones muertas y bloat}

Vacuum no elimina una versión que aún puede ser visible. Una transacción \emph{idle in transaction} puede ser más dañina que una consulta que termina, porque conserva un snapshot sin producir trabajo útil.

\textbf{Bloat} es el crecimiento físico innecesario de tablas e índices cuando versiones que ya no representan el estado vigente continúan ocupando páginas. No toda versión antigua es bloat: si un snapshot activo todavía puede verla, es historia necesaria. El problema aparece cuando versiones ya reclamables se acumulan o cuando snapshots antiguos retrasan su limpieza.

Explíquelo con una frase:

\begin{quote}
``MVCC cambia parte del costo de bloquear por el costo de conservar y limpiar historia; bloat es historia que ya no resulta útil, pero todavía ocupa espacio y genera trabajo.''
\end{quote}

Consecuencias: más páginas leídas, menor eficiencia de caché, índices mayores, más E/S, backups más pesados y mayor trabajo de mantenimiento. Un \texttt{VACUUM} normal suele volver espacio reutilizable dentro de la relación, pero no necesariamente devuelve inmediatamente el tamaño del archivo al sistema operativo.

\subsubsection*{Serializabilidad}

Serializabilidad es una propiedad global de la ejecución. Puede implementarse mediante locks de dos fases, validación, SSI u otros mecanismos. MVCC aporta versiones, pero la garantía depende de cómo se detectan dependencias peligrosas.

\subsection{Índices y EXPLAIN en profundidad}

\subsubsection*{Árbol B, fan-out y relación con MVCC}

Un árbol B es un árbol balanceado de múltiples hijos. Los nodos internos conservan claves separadoras y referencias; las hojas conservan claves ordenadas y referencias a filas o páginas. El alto fan-out mantiene poca altura incluso con millones de entradas.

No confunda localización con visibilidad. El índice produce candidatos; la tabla o el índice, según el plan, aportan datos; las reglas MVCC deciden qué versiones puede observar el snapshot. Un index scan puede tener que avanzar más de veinte entradas para reunir veinte filas visibles.

Para CampusShop, el orden \texttt{(student\_id, status, confirmed\_at DESC)} se deriva del patrón exacto: dos igualdades, un orden y un límite. No es una regla universal.

\subsubsection*{Qué estima el planificador}

El planificador utiliza estadísticas de cardinalidad, distribución, selectividad y costos configurados. Un error de estimación puede elegir un plan pobre aun cuando la estructura exista.

\subsubsection*{Por qué leer de abajo hacia arriba}

Los nodos inferiores producen filas para los superiores. Debe preguntar en cada nodo:

\begin{itemize}
\item ¿cuántas filas recibe?;
\item ¿cuántas produce?;
\item ¿cuántas veces se ejecuta?;
\item ¿qué buffers toca?;
\item ¿dónde aparece el mayor descarte o multiplicación?
\end{itemize}

\subsubsection*{EXPLAIN ANALYZE puede modificar datos}

Si se ejecuta sobre \texttt{INSERT}, \texttt{UPDATE} o \texttt{DELETE}, la sentencia realmente se ejecuta. Para una demo segura use SELECT o una transacción que luego se deshace explícitamente.

\subsubsection*{Seq Scan no es fracaso}

Puede ser la mejor opción para tablas pequeñas o consultas poco selectivas. El objetivo no es ``forzar Index Scan'', sino minimizar costo real.

\subsubsection*{Índice compuesto}

La utilidad depende del prefijo de columnas, igualdad, rango, orden y cobertura. El índice de CampusShop se justifica por el patrón exacto de filtros y \texttt{ORDER BY}; no por una regla general de poner todas las columnas de la consulta.

\subsubsection*{Otros scans}

Ante preguntas avanzadas, mencione:

\begin{itemize}
\item index-only scan: puede responder desde el índice si visibilidad y columnas lo permiten;
\item bitmap scan: combina muchas coincidencias del índice antes de visitar páginas;
\item heap/table lookup: un índice puede localizar filas, pero todavía requerir acceso a la tabla.
\end{itemize}

No es necesario desarrollar estos scans en la ruta esencial.

\subsection{LSM, Bloom y compactación en profundidad}

\subsubsection*{¿Una LSM-tree es realmente un árbol?}

El nombre histórico describe una familia de estructuras en niveles. En implementaciones modernas, el estado persistente suele verse como conjuntos de SSTables por niveles más estructuras auxiliares, no como un árbol de nodos al estilo B-tree.

\subsubsection*{SSTable e inmutabilidad}

Una SSTable no se modifica in-place. Nuevas versiones se escriben en archivos posteriores; compactación produce archivos nuevos y retira los antiguos cuando es seguro.

\subsubsection*{Merge ordenado}

Si cada entrada está ordenada por clave, el merge compara cabeceras y avanza linealmente. Con $k$ runs puede usarse una cola de prioridad para seleccionar la menor clave. Esto es distinto de ordenar nuevamente todos los datos.

\subsubsection*{Filtro de Bloom}

Propiedades que debe poder explicar:

\begin{itemize}
\item estructura probabilística de membresía;
\item utiliza varios hashes y un vector de bits;
\item ausencia definitiva cuando algún bit requerido es cero;
\item presencia posible cuando todos son uno;
\item falso positivo posible;
\item falso negativo no admisible en uso correcto;
\item utilidad principal en búsquedas exactas de claves ausentes;
\item no devuelve el valor ni resuelve visibilidad.
\end{itemize}

La tasa de falsos positivos depende de bits por clave y cantidad de funciones hash. No hace falta derivar la fórmula durante la clase, pero puede mencionar que más memoria reduce lecturas innecesarias.

\subsubsection*{Tombstone y resurrección}

Un tombstone debe alcanzar niveles donde existan versiones antiguas antes de eliminarse. Si se descarta demasiado pronto, una compactación futura puede dejar visible un valor que se pretendía borrar.

\subsubsection*{Compaction leveled y tiered}

Leveled limita solapamiento y favorece lecturas/espacio, a costa de más reescritura. Tiered agrupa runs con menos reescritura temprana, pero puede aumentar archivos consultados y espacio temporal. No presentar una como superior universal.

\subsubsection*{Stall y estabilidad}

Un stall es similar a backpressure en una cola. Si la tasa de entrada $\lambda$ supera sostenidamente la tasa de servicio $\mu$, la cola crece. El motor debe reducir $\lambda$ efectiva, aumentar $\mu$ o aceptar inestabilidad.

La pendiente del backlog es una métrica causal:

\[
\frac{dB}{dt} \approx \lambda - \mu.
\]

Un backlog grande pero decreciente puede ser menos preocupante que uno menor con pendiente positiva persistente.

\subsection{Cómo explicar los conceptos en tres profundidades}

\subsubsection*{WAL}

\textbf{30 segundos.} ``El motor anota cómo recuperar un cambio antes de depender de la página final; confirma al cliente solo cuando el COMMIT está durable.''

\textbf{2 minutos.} Agregar las dos reglas, prefijo durable, páginas atrasadas y REDO.

\textbf{5 minutos.} Construir timeline, mover crash, explicar \texttt{flushLSN}, \texttt{commitLSN}, pageLSN, group commit e idempotencia 9F2A.

\subsubsection*{MVCC}

\textbf{30 segundos.} ``Conserva versiones para que una lectura mantenga una historia coherente mientras otras transacciones escriben.''

\textbf{2 minutos.} Definir snapshot, visibilidad, lector antiguo/nuevo y costo de retener versiones.

\textbf{5 minutos.} Agregar niveles de aislamiento, write skew, VACUUM, bloat y locks que todavía existen.

\subsubsection*{EXPLAIN e índice}

\textbf{30 segundos.} ``El índice agrega una ruta; EXPLAIN muestra si el planificador piensa usarla y ANALYZE mide la ejecución.''

\textbf{2 minutos.} Leer scan, filtro, sort, limit, filas y buffers.

\textbf{5 minutos.} Agregar selectividad, estadísticas, índice compuesto, seq scan válido y costo de escritura.

\subsubsection*{LSM, Bloom y compactación}

\textbf{30 segundos.} ``La LSM escribe rápido en memoria y archivos inmutables; luego debe fusionarlos. Bloom evita abrir algunos archivos cuando la clave está ausente.''

\textbf{2 minutos.} Recorrer write path, SSTables, read path, tombstones y falsos positivos.

\textbf{5 minutos.} Agregar amplificaciones, leveled/tiered, backlog slope, archivos L0, stall y SLO.

\section{Preguntas socráticas}

Las siguientes preguntas se utilizan para conducir la clase y recuperar el razonamiento del estudiante antes de ofrecer una respuesta técnica. Conviene formularlas sobre el caso CampusShop, esperar una hipótesis y recién después completar la explicación.

\subsection{Banco de preguntas previsibles}

\begin{enumerate}
\item \textbf{¿COMMIT significa que la página está en SSD?} No. Significa que el motor decidió confirmar y, para durabilidad, que la evidencia requerida está persistida. Las páginas pueden escribirse después.

\item \textbf{¿Por qué no escribir directamente las páginas?} Es posible, pero sincronizar muchas páginas dispersas en cada commit suele ser más costoso y complica atomicidad. WAL concentra evidencia secuencial.

\item \textbf{¿WAL es un backup?} No. WAL ayuda a recovery y puede participar en recuperación a un punto, pero requiere una base consistente y políticas de retención; no sustituye backups independientes.

\item \textbf{¿WAL garantiza exactamente una vez?} No. Recupera la historia del motor. La identidad 9F2A y su persistencia atómica sostienen exactamente-una-vez lógica.

\item \textbf{¿Qué pasa si se pierde el ACK?} La transacción puede estar confirmada. El cliente reintenta y el servidor devuelve el resultado asociado a 9F2A.

\item \textbf{¿El flush durable empieza en el UPDATE?} No. El update se agrega antes; el flush posterior avanza la frontera durable hasta el COMMIT e incluye el prefijo previo.

\item \textbf{¿Qué pasa si el crash ocurre durante recovery?} Recovery debe poder repetirse; pageLSN y registros idempotentes evitan aplicar indebidamente el mismo cambio.

\item \textbf{¿Checkpoint implica RPO cero?} No. Reduce trabajo de recuperación. El RPO depende de qué commits eran durables y de la falla considerada.

\item \textbf{¿MVCC elimina locks?} No. Reduce contención lector--escritor, pero siguen existiendo locks de escritura, DDL e invariantes.

\item \textbf{¿Snapshot es una copia?} No. Es una regla de visibilidad.

\item \textbf{¿MVCC es serializable?} No necesariamente. Snapshot Isolation puede admitir write skew.

\item \textbf{¿Por qué VACUUM no borra todo enseguida?} Porque una transacción activa puede necesitar una versión anterior.

\item \textbf{¿Un índice siempre acelera?} No. Puede no ser selectivo, la tabla puede ser pequeña y mantenerlo siempre cuesta.

\item \textbf{¿EXPLAIN ANALYZE es seguro?} En una consulta \texttt{SELECT}, normalmente sí con los cuidados habituales. En escrituras, ejecuta el cambio; debe usarse en una transacción o entorno controlado.

\item \textbf{¿Por qué elegir un escaneo secuencial si hay índice?} Porque leer secuencialmente gran parte de una tabla puede ser más barato que alternar entre índice y páginas.

\item \textbf{¿LSM elimina escrituras aleatorias?} Reduce parte del trabajo aleatorio del camino crítico, pero la compactación reescribe datos y puede producir amplificación.

\item \textbf{¿Bloom puede afirmar ausencia cuando la clave sí está?} No debería; eso sería un falso negativo. Puede responder ``posiblemente presente'' cuando la clave no está.

\item \textbf{¿Qué significa omitir una SSTable?} No abrirla para esa búsqueda concreta porque Bloom demostró ausencia. No significa borrar el archivo.

\item \textbf{¿Por qué no borrar un tombstone inmediatamente?} Porque una versión antigua puede sobrevivir en otro nivel y resucitar.

\item \textbf{¿Un stall es una falla?} Es una protección. Se vuelve síntoma de incapacidad cuando domina el SLO o aparece continuamente.

\item \textbf{Si las páginas no se materializaron, ¿cómo aparece toda la data después del crash?} Recovery parte de páginas persistentes anteriores y aplica registros WAL de modificación confirmados. COMMIT marca la decisión; no contiene por sí solo toda la fila.

\item \textbf{¿Qué es bloat?} Acumulación física de versiones muertas, entradas de índice obsoletas y espacio poco aprovechado. MVCC necesita historia; bloat es historia que ya no es necesaria y todavía cuesta.

\item \textbf{¿El índice decide qué versión devuelve MVCC?} No. El índice localiza candidatos; la regla de visibilidad decide cuál puede observar la transacción.

\item \textbf{¿Por qué el árbol B tiene pocos niveles?} Porque cada página interna tiene alto fan-out y puede dirigir la búsqueda hacia muchos hijos.

\item \textbf{¿Cuál es mejor: B-tree o LSM?} Depende de carga, consulta, hardware, mantenimiento y objetivos. La pregunta correcta es qué costo conviene desplazar.
\end{enumerate}

\section{Parte C: evaluación, auditoría y evidencias}

\subsection*{Parte C: auditoría}

\section{Criterios de evaluación formativa}

La evaluación formativa observa si el estudiante relaciona una promesa con la falla considerada, el mecanismo que la sostiene, el costo desplazado y la evidencia que permitiría verificarla.

\subsection{Alineación constructiva}

\begin{center}
\small
\begin{tabular}{p{0.20\textwidth}p{0.25\textwidth}p{0.20\textwidth}p{0.18\textwidth}}
\toprule
Resultado & Actividad & Evidencia & Criterio \\
\midrule
Distinguir capas & Clasificar decisiones distribuidas/locales & Diagrama anotado & Autoridad y persistencia separadas \\
Explicar WAL & Completar timeline y mover crash & Prefijo durable y ACK ubicados & Todas y solo confirmadas \\
Explicar MVCC & Timeline lector/escritor & Versiones visibles justificadas & Regla de snapshot, no slogan \\
Leer planes & Comparar Seq/Index Scan & Filas, buffers, sort y costo & Evidencia antes que tiempo aislado \\
Diagnosticar LSM & Leer dashboard y calcular pendiente & Hipótesis y métricas & Cadena causal defendible \\
Diseñar & CampusShop vs telemetría & Defensa de 90 segundos & Falla, garantía, costo y métrica \\
\bottomrule
\end{tabular}
\end{center}

\subsection{Rúbrica de logro}

\begin{center}
\begin{tabular}{p{0.18\textwidth}p{0.68\textwidth}}
\toprule
Nivel & Evidencia \\
\midrule
Inicial & Nombra componentes o productos sin garantía ni falla. \\
En desarrollo & Conecta mecanismo con problema, pero no ubica frontera o costo. \\
Logrado & Explica mecanismo, falla, garantía, límite y métrica. \\
Avanzado & Compara alternativas, cuantifica costo y propone experimento refutable. \\
\bottomrule
\end{tabular}
\end{center}

\section{Auditoría y evidencias}

\subsection{Preguntas que puede formular un auditor pedagógico}

\begin{enumerate}
\item \textbf{¿Por qué comienza con un caso y no con una definición?} Porque el caso activa conocimientos previos y crea una necesidad cognitiva. La formalización responde después a una pregunta ya comprendida.

\item \textbf{¿Cómo verifica aprendizaje durante la clase?} Mediante dos timelines completados, lectura de planes, cálculo de backlog, defensa grupal y exit ticket.

\item \textbf{¿Cómo evita que la analogía simplifique demasiado?} Cada intuición se acompaña de una definición técnica y un límite explícito.

\item \textbf{¿Por qué incluye productos reales?} PostgreSQL y RocksDB muestran instancias concretas de principios transferibles; no se evalúa memorizar comandos.

\item \textbf{¿Cómo atiende distintos niveles previos?} El vocabulario explícito reduce barreras, mientras las diapositivas de extensión y el dossier permiten profundizar sin bloquear la ruta esencial.

\item \textbf{¿Cómo se conecta con clases anteriores?} Cada una aporta una promesa que necesita persistencia: incertidumbre, idempotencia, replicación, consenso y ACID.

\item \textbf{¿Qué evidencia conservará?} Prediagnóstico, timelines, planes anotados, cálculo grupal, defensa y exit ticket.

\item \textbf{¿Cómo se relaciona con el Trabajo Final?} Los estudiantes deben declarar carga, falla, garantía de durabilidad, estrategia de datos, recovery, métricas y condición de revisión.
\end{enumerate}

\subsection{Preguntas técnicas que puede formular un auditor}

\begin{enumerate}
\item Explique por qué un COMMIT durable puede coexistir con páginas atrasadas.
\item Distinga flush común de flush durable.
\item Explique por qué el flush del COMMIT hace durable un prefijo del WAL.
\item ¿Qué ocurre entre append de COMMIT y flush si hay crash?
\item ¿Cómo se recupera un ACK perdido sin duplicar compra?
\item ¿Qué diferencia hay entre checkpoint y backup?
\item ¿Qué locks siguen existiendo con MVCC?
\item ¿Por qué Snapshot Isolation no es serializable?
\item ¿Cómo detectaría estadísticas erróneas en \texttt{EXPLAIN ANALYZE}?
\item ¿Cuándo un seq scan es mejor?
\item ¿Qué falso resultado puede dar Bloom?
\item ¿Por qué un tombstone debe propagarse?
\item ¿Qué revela la pendiente del backlog?
\item ¿Por qué un stall puede ser sano?
\item ¿Cómo cambia el diseño si la falla es pérdida de región y no crash local?
\item Si solo el registro COMMIT fuera durable y faltaran los registros de modificación, ¿podría recuperarse la transacción? Justifique.
\item Distinga versiones antiguas necesarias de bloat reclamable.
\item Explique el recorrido consulta--planificador--árbol B--candidatos--visibilidad MVCC.
\end{enumerate}

Las respuestas están desarrolladas en la Parte B. Antes de la clase, practique responder cada una en menos de dos minutos.

\subsection{Evidencias que conviene conservar}

\begin{itemize}
\item fotografía del timeline WAL con la frontera durable marcada;
\item dos respuestas distintas al caso de crash y su corrección;
\item timeline MVCC con versiones visibles;
\item comparación de plan antes/después del índice;
\item cálculo escrito de 30 MB/s de déficit y tiempos a umbrales;
\item defensa CampusShop/telemetría;
\item resumen anónimo del exit ticket;
\item ajuste para la clase 7 basado en una duda recurrente.
\end{itemize}

\subsection{Matriz de corrección y preguntas de defensa}

La corrección debe evaluar razonamiento y no solo vocabulario. Use la siguiente matriz:

\begin{center}
\small
\begin{tabular}{p{0.23\textwidth}p{0.29\textwidth}p{0.31\textwidth}}
\toprule
Dimensión & Evidencia mínima & Pregunta de defensa \\
\midrule
Problema y objeto & Identifica transacción, página, versión, consulta o SSTable & ¿Qué objeto exacto cambia o se recupera? \\
Falla & Declara crash, pérdida de energía, timeout o sobrecarga & ¿En qué punto de la ejecución coloca la falla? \\
Garantía & Formula el comportamiento observable correcto & ¿Qué debe seguir siendo verdadero después de la falla? \\
Mecanismo & Explica orden, frontera o regla de visibilidad & ¿Qué parte del mecanismo sostiene esa garantía? \\
Costo & Señala espera, I/O, espacio, mantenimiento o complejidad & ¿Dónde se paga la decisión? \\
Evidencia & Propone métrica, plan, timeline o experimento & ¿Qué observación podría refutar su hipótesis? \\
\bottomrule
\end{tabular}
\end{center}

Para cada ejercicio, una respuesta lograda debe contener:

\begin{enumerate}
\item objeto exacto: transacción, página, versión, consulta o SSTable;
\item falla y alcance;
\item garantía observable;
\item mecanismo con orden o frontera;
\item costo desplazado;
\item evidencia o experimento.
\end{enumerate}

Errores graves que deben corregirse explícitamente:

\begin{itemize}
\item ``COMMIT significa página en SSD'';
\item ``WAL da exactamente una vez'';
\item ``flush durable empieza en el UPDATE'';
\item ``MVCC no usa locks'';
\item ``snapshot es copia de la base'';
\item ``índice siempre mejora'';
\item ``EXPLAIN ANALYZE solo estima'';
\item ``Bloom puede descartar una clave existente'';
\item ``archivo imposible'' sin aclaración;
\item ``stall aumenta capacidad'';
\item ``LSM siempre es más rápida''.
\end{itemize}

\subsection{Errores previsibles}

Los errores siguientes son diagnósticos del modelo mental del estudiante. No basta con indicar que la respuesta es incorrecta; conviene volver al caso, mover la falla y pedir evidencia.

\begin{itemize}
\item \textbf{``COMMIT significa página en SSD''.} Separar registro de commit durable de materialización posterior de páginas.
\item \textbf{``Replicar en RAM vuelve durable''.} Distinguir cantidad de copias, persistencia y falla común.
\item \textbf{``WAL garantiza exactamente una vez''.} Volver a la identidad 9F2A y a la atomicidad entre efecto y registro idempotente.
\item \textbf{``MVCC elimina todos los locks y anomalías''.} Recuperar write skew, DDL e invariantes.
\item \textbf{``Un índice siempre mejora''.} Pedir el plan, selectividad y costo sobre inserciones.
\item \textbf{``Bloom confirma que la clave existe''.} Recordar que solo puede demostrar algunas ausencias; un positivo significa ``posiblemente presente''.
\item \textbf{``LSM elimina el costo de escribir''.} Mostrar compaction, amplificación y backlog.
\item \textbf{``Un stall es una falla del motor''.} Explicar que es backpressure protector; su frecuencia puede revelar falta de capacidad.
\end{itemize}

\subsection{Contingencias}

\subsubsection*{Si falta tiempo}

Eliminar detalles de STEAL/NO-STEAL, scans avanzados y políticas de compactación. No eliminar vocabulario, timelines, Bloom, lectura de plan ni cierre.

\subsubsection*{Si falla PostgreSQL}

Usar capturas o salidas impresas. Pedir primero una predicción sobre plan y luego revelar la salida.

\subsubsection*{Si el grupo participa poco}

Usar voto individual anónimo, discusión en parejas y selección de dos respuestas escritas. No depender de voluntarios espontáneos.

\subsubsection*{Si el grupo ya conoce las siglas}

Acelerar definiciones, pero no omitir límites. Pedir contraejemplos: MVCC sin serializabilidad, índice ignorado, Bloom falso positivo y WAL sin idempotencia.

\subsubsection*{Si surgen preguntas demasiado específicas de producto}

Responder el principio y ubicar el detalle como extensión. Ejemplo: ante una pregunta de configuración de RocksDB, volver a tasa de ingreso, capacidad, backlog, amplificación y SLO.

\subsection{Checklist personal de estudio}

Antes de dictar, confirmar que puede explicar sin mirar:

\begin{itemize}
\item la diferencia entre storage distribuido y motor local;
\item el significado de todas las siglas;
\item la identidad 9F2A;
\item COMMIT, log durable y flush durable;
\item las dos reglas de WAL;
\item el timeline completo y cuatro crashes;
\item la relación \texttt{flushLSN}--\texttt{commitLSN};
\item group commit y checkpoint;
\item snapshot, visibilidad y serializabilidad;
\item VACUUM, bloat y write skew;
\item lectura de ambos planes;
\item DDL del índice compuesto;
\item write/read/space amplification;
\item ruta de escritura y lectura LSM;
\item merge ordenado, tombstone y filtro de Bloom;
\item cálculo de backlog y umbrales;
\item leveled frente a tiered;
\item una defensa completa para CampusShop y otra para telemetría.
\end{itemize}

\subsection{Trabajo Final Integrador}

La Clase 06 debe dejar una evidencia reutilizable en el Trabajo Final Integrador. Cada grupo debe incorporar, para al menos un componente de su sistema:

\begin{enumerate}
\item patrón de lectura y escritura, volumen, crecimiento y skew;
\item dato crítico e invariante asociado;
\item frontera de durabilidad y modelo de falla;
\item estrategia de recuperación y pérdida máxima tolerable;
\item elección razonada entre árbol B, LSM o una combinación;
\item costo desplazado: WAL, índices, versiones, compactación o stalls;
\item dos métricas y un experimento de falla o carga;
\item condición explícita para revisar la decisión.
\end{enumerate}

No se acepta ``usaremos X porque escala''. En la defensa, el grupo debe responder qué garantiza X, ante qué falla, dónde se paga y qué evidencia mostraría que la elección fue incorrecta.

\subsection{Bibliografía priorizada para el docente}

Orden sugerido de estudio:

\begin{enumerate}
\item Kleppmann, \emph{Designing Data-Intensive Applications}, capítulo 3.
\item PostgreSQL Documentation, \emph{Write-Ahead Logging}.
\item PostgreSQL Documentation, \emph{MVCC} y \emph{Routine Vacuuming}.
\item Kleppmann, capítulo 7, para aislamiento y write skew.
\item O'Neil et al., \emph{The Log-Structured Merge-Tree}.
\item RocksDB Wiki: compaction, compaction stats y write stalls.
\item Gray y Reuter, secciones de logging y recovery, para preguntas avanzadas.
\end{enumerate}

\subsection{Exit ticket y seguimiento}

Preguntas:

\begin{enumerate}
\item diferencia entre aceptado, escrito y durable;
\item costo oculto de MVCC y una métrica;
\item diferencia entre \texttt{EXPLAIN} y \texttt{EXPLAIN ANALYZE};
\item señales de backlog de compactación.
\end{enumerate}

Clasifique las respuestas en:

\begin{itemize}
\item concepto claro;
\item vocabulario correcto, razonamiento incompleto;
\item confusión de frontera;
\item error técnico que debe retomarse.
\end{itemize}

Comience la clase 7 devolviendo una duda real, por ejemplo:

\begin{quote}
``Varios distinguieron WAL durable de página escrita, pero todavía quedó una pregunta: ¿qué ocurre cuando esa historia deja de ser local y la consumen otros procesos?''
\end{quote}

\begin{uakeybox}
\textbf{Criterio final de éxito docente.} La clase salió bien si los estudiantes no solo pueden definir WAL o MVCC, sino también mover una falla por una ejecución, justificar qué debe sobrevivir, leer evidencia y defender una decisión con límites explícitos.
\end{uakeybox}

\end{document}
